% §8 Discussion
\subsection{Practitioner Guidance}

Our analysis yields concrete guidance for teams building cascade voice agents.

\paragraph{When cascade is sufficient.}
A cascade agent without tool use ($k=0$) with streaming TTS can achieve a silence gap as low as $\Lambda^* \approx 280$~ms (using representative values from Table~\ref{tab:latency_example} with $\lambda_{tts}^{irr}$ reduced to 0 by full streaming).
This is close to but not within the 200~ms threshold.
For applications where users tolerate slightly longer gaps (e.g., task-oriented dialogue where cognitive load reduces gap sensitivity), a streaming-TTS cascade is a reasonable choice.

\paragraph{When tool use is required.}
For $k \geq 1$, the silence gap grows by at least $\lambda_{lm}^{irr} \approx 100$~ms per step.
Filler generation is strongly recommended for any tool-using cascade agent to mask the silence gap perceptually.
Parallel tool execution should be implemented wherever tool calls are independent.
For $k \geq 3$, the irreducible floor exceeds 500~ms regardless of optimisations.
Users will perceive the delay as significant.
In these cases, hybrid routing to a native voice model for tool-intensive turns is the only structural solution.

\paragraph{When to choose native voice.}
Native voice models are justified in three cases.
First, the CCC cannot be satisfied by any combination of the structural levers (streaming TTS and hybrid routing).
Second, the turn mix is dominated by short, conversational turns where the overhead of the cascade pipeline is felt most acutely.
Third, barge-in and interruption handling quality is a hard requirement, since native models handle these natively without a separate VAD pipeline.
The higher cost and reduced auditability of native models are the primary tradeoffs.

\subsection{Limitations}

The perceptual threshold $\tau_{human} = 200$~ms is derived from studies of natural human-human conversation \citep{levinson2015timing,stivers2009universals}.
Human-agent interaction may have different norms.
Users may calibrate expectations downward when they know they are speaking to an AI.
Our analysis treats $\tau_{human}$ as a constant, but future work should treat it as a distribution conditioned on user, context, and task type.

The latency decomposition in Definition~\ref{def:decomposition} assumes that irreducible latencies are stable across inputs.
In practice, $\lambda_{lm}^{irr}$ grows with context length, and $\lambda_{asr}^{irr}$ grows with utterance complexity.
A more complete model would express irreducible latencies as functions of input characteristics.

Our model does not account for network latency between pipeline components, which adds a reducible but non-negligible contribution in cloud-hosted deployments.
